\section{Make decompression an exact-size call}
Call \texttt{png\_inflate} with the joined zlib stream address in \texttt{\$a0}, its
unsigned byte count in \texttt{\$a1}, the row-workspace address in \texttt{\$a2}, and
the exact expected output count in \texttt{\$a3}. For the worked image, the two counts
are 46 and 35. They have different sources: the collector sums IDAT payload lengths,
whereas the dimensions predict $h(1+3w)$ inflated bytes. Passing the compressed count
as the expected row size would request eleven bytes that this image does not produce.

A successful call returns zero in \texttt{\$v0} and the exact produced count in
\texttt{\$v1}. Failure returns one and zero. Preserve the stream, saved registers and
stack pointer. Output bytes inside the selected row workspace may change on failure;
the caller must not reconstruct or display them after a nonzero status. Leave bytes
outside that workspace unchanged. The stream, output and private stack storage must
not overlap. Counts and addresses retain all 64 bits. There are no file operations
or framebuffer writes in this entry.

Require between six bytes and 1 MiB of input and at most 786944 expected output bytes.
Reject a count outside these bounds before reading the stream. A zero expected count
is useful for testing a valid compressed empty stream, although a supported PNG with
positive dimensions always predicts some row bytes. The caller supplies readable input
and writable output windows of the selected sizes. Bounds on their counts do not make
an arbitrary pointer readable or allocate its destination automatically.

Checksum agreement cannot replace a produced-byte count. One test stream expands to
one ASCII A byte, while the caller requests 65522 bytes and initializes the workspace
to decimal 165. Hashing that A followed by 65521 unchanged bytes of value 165 gives the
same Adler-32 as hashing the A alone. The extra bytes were never produced by the
stream. Compare the actual output cursor with the expected end before accepting the
checksum; reading an untouched part of the workspace does not turn it into output.

Use each call's own workspace for code lengths, counts and ordered symbols. A previous
fixed block's table must not supply values for a later dynamic block, and a failed call
must not leave parser state that changes a later call. Read-only base/extra-bit tables
can be shared. A particular stack layout is an implementation choice; the interface
specifies the source, output and returned state, rather than naming the reference
implementation's private slots.
